\chapter{Hyperemesis Gravidarum}
\index{Hyperemesis Gravidarum}

\section*{Definition and Overview}
Hyperemesis gravidarum (HG) is a severe form of nausea and vomiting in pregnancy that goes far beyond ordinary morning sickness. It causes persistent, severe vomiting that leads to dehydration, weight loss, and electrolyte imbalance, and it can seriously affect a woman's health and quality of life. HG affects around 1--2\% of pregnancies and is the most common reason for hospital admission in early pregnancy. Treatment includes anti-nausea medicines, fluids, and nutritional support, and with good management most women improve as the pregnancy progresses.

\section*{Epidemiology}
Around 70--80\% of pregnant women experience nausea, and up to 50\% vomit during pregnancy. Hyperemesis gravidarum, the severe form, affects approximately 1--2\% of pregnancies. It usually begins between 4 and 8 weeks of pregnancy and typically improves by 16--20 weeks, though some women have symptoms throughout pregnancy. HG is more common in first pregnancies, multiple pregnancies, and in women with a family history. It recurs in around 80\% of subsequent pregnancies.

\section*{Aetiology and Pathophysiology}
\begin{itemize}
    \item \textbf{hCG:} high levels of pregnancy hormone hCG are strongly associated with HG
    \item \textbf{Oestrogen:} high oestrogen levels may contribute
    \item \textbf{Thyroid:} transient thyroid changes occur in many women with HG
    \item \textbf{Gastric motility:} pregnancy hormones slow stomach emptying
    \item \textbf{Genetic factors:} HG runs in families, suggesting a genetic component
    \item \textbf{Psychological factors:} stress and anxiety can worsen symptoms, but they are not the cause
    \item \textbf{Pathophysiology:} the exact mechanism is unclear, but hormonal and genetic factors drive severe vomiting and its consequences
\end{itemize}

\section*{Clinical Features}
\begin{itemize}
    \item Persistent, severe nausea and vomiting, often worse than in normal pregnancy
    \item Inability to keep down food or fluids
    \item Weight loss of more than 5\% of pre-pregnancy weight
    \item Signs of dehydration: dark urine, dry mouth, and dizziness
    \item Electrolyte imbalance and weakness
    \item Reduced urine output
    \item Loss of saliva and drooling in some women
    \item Symptoms that seriously affect daily life and mood
\end{itemize}

\section*{Differential Diagnosis}
\begin{itemize}
    \item Ordinary morning sickness, which is milder and does not cause dehydration
    \item Gastroenteritis and other infections
    \item Thyroid disease, which can mimic or accompany HG
    \item Gastrointestinal conditions, including reflux and ulcers
    \item Molar pregnancy, which causes high hCG and severe symptoms
    \item Investigation excludes other causes before HG is diagnosed
\end{itemize}

\section*{When to Seek Medical Care}
Pregnant women who cannot keep down fluids, are losing weight, or feel unwell should seek medical review early, as HG responds better to early treatment.

\textbf{Contact your local emergency services for:}
\begin{itemize}
    \item Inability to keep down any fluids with signs of severe dehydration
    \item Passing very little or no urine
    \item Fainting or severe dizziness
    \item Confusion or altered consciousness
    \item Abdominal pain, fever, or blood in vomit
    \item Thoughts of harming yourself; HG is distressing and support is available
\end{itemize}

\section*{Diagnosis and Investigations}
\begin{itemize}
    \item \textbf{Clinical assessment:} history of vomiting, weight, and hydration
    \item \textbf{Urine tests:} ketones indicate starvation and dehydration
    \item \textbf{Blood tests:} electrolytes, kidney function, and liver tests
    \item \textbf{Thyroid tests:} to exclude hyperthyroidism
    \item \textbf{Ultrasound:} confirms the pregnancy and excludes molar pregnancy
    \item \textbf{Psychological assessment:} where mood is significantly affected
\end{itemize}

\section*{Management}
\begin{itemize}
    \item \textbf{Anti-nausea medicines:} safe options include ondansetron, metoclopramide, and doxylamine, chosen with pregnancy safety in mind
    \item \textbf{Fluids:} oral rehydration where possible; intravenous fluids for dehydration
    \item \textbf{Vitamins:} thiamine and other B vitamins to prevent deficiency
    \item \textbf{Nutrition:} small frequent meals and foods that are tolerated
    \item \textbf{Hospital care:} intravenous fluids, electrolytes, and monitoring for severe cases
    \item \textbf{Psychological support:} counselling for the distress HG causes
    \item \textbf{Alternative approaches:} ginger and acupressure help some women
    \item \textbf{Ongoing care:} regular review throughout the pregnancy
\end{itemize}

\section*{Complications}
\begin{itemize}
    \item Dehydration, electrolyte imbalance, and kidney injury
    \item Weight loss and malnutrition
    \item Vitamin deficiencies, including thiamine deficiency
    \item Oesophageal tears and bleeding from repeated vomiting
    \item Blood clots from dehydration and immobility
    \item Depression, anxiety, and post-traumatic stress
    \item Impact on work, relationships, and quality of life
    \item Rarely, effects on the baby, reduced by early treatment
\end{itemize}

\section*{Prognosis and Outcomes}
Most women with HG improve by 16--20 weeks, and with early, aggressive treatment, the risks to mother and baby are low. Babies born to women with well-managed HG have outcomes similar to other pregnancies. A minority of women have symptoms throughout pregnancy, and recurrence in future pregnancies is common, allowing early planning. Psychological support is important, as HG is a distressing condition even when it resolves.

\section*{Prevention and Self-Care}
\begin{itemize}
    \item If HG occurred in a previous pregnancy, plan early treatment with your doctor
    \item Start anti-nausea treatment early
    \item Eat small, frequent meals and avoid triggers
    \item Sip fluids regularly; try ice chips or rehydration drinks
    \item Rest and avoid strong smells
    \item Seek medical review early rather than struggling
    \item Ask for help with daily tasks and emotional support
    \item Carry a written plan for flare-ups
\end{itemize}

\section*{Sources and Further Reading}
The following reputable sources provide further information for healthcare students and patients seeking reliable, current advice about hyperemesis gravidarum.

\begin{itemize}
    \item Healthdirect. \textit{Hyperemesis gravidarum.} https://www.healthdirect.gov.au/hyperemesis-gravidarum
    \item The Royal Women's Hospital. \textit{Hyperemesis gravidarum fact sheet.} https://www.thewomens.org.au/
    \item Pregnancy, Birth and Baby. \textit{Nausea and vomiting in pregnancy.} https://www.pregnancybirthbaby.org.au/
    \item HER Foundation. \textit{Hyperemesis gravidarum information and support.} https://www.hyperemesis.org/
    \item NHS. \textit{Severe vomiting in pregnancy.} https://www.nhs.uk/pregnancy/
    \item Mayo Clinic. \textit{Hyperemesis gravidarum: symptoms and causes.} https://www.mayoclinic.org/
\end{itemize}
